
\documentclass[18pt,aspectratio=169]{beamer}
%\usetheme{kit}
\usepackage{ngerman}

% metropolis theme: https://github.com/matze/mtheme
\usetheme[sectionpage=progressbar,subsectionpage=progressbar,progressbar=frametitle]{metropolis}

% Notizen
%\setbeameroption{show notes}
%\setbeameroption{show notes on second screen=right}

\setbeamertemplate{section in toc}[sections numbered]
\setbeamertemplate{subsection in toc}[subsections numbered]

%redefine note page with bigger slide
\makeatletter
\defbeamertemplate*{note page}{mynotes}
{%
	{%
		\scriptsize
		\usebeamerfont{note title}\usebeamercolor[fg]{note title}%
		\ifbeamercolorempty[bg]{note title}{}{%
			\insertvrule{.45\paperheight}{note title.bg}%
			\vskip-.45\paperheight%
			\nointerlineskip%
		}%
		\vbox{
			\hfill\insertslideintonotes{0.45}\hskip-\Gm@rmargin\hskip0pt%
			\vskip-0.45\paperheight%
			\nointerlineskip
			\begin{pgfpicture}{0cm}{0cm}{0cm}{0cm}
				\begin{pgflowlevelscope}{\pgftransformrotate{90}}
					{\pgftransformshift{\pgfpoint{-2cm}{0.2cm}}%
						\pgftext[base,left]{\usebeamerfont{note date}\usebeamercolor[fg]{note date}\the\year-\ifnum\month<10\relax0\fi\the\month-\ifnum\day<10\relax0\fi\the\day}}
				\end{pgflowlevelscope}
		\end{pgfpicture}}
		\nointerlineskip
		\vbox to .45\paperheight{\vskip0.5em
			\hbox{\insertshorttitle[width=8cm]}%
			\setbox\beamer@tempbox=\hbox{\insertsection}%
			\hbox{\ifdim\wd\beamer@tempbox>1pt{\hskip4pt\raise3pt\hbox{\vrule
						width0.4pt height7pt\vrule width 9pt
						height0.4pt}}\hskip1pt\hbox{\begin{minipage}[t]{7.5cm}\def\breakhere{}\insertsection\end{minipage}}\fi%
			}%
			\setbox\beamer@tempbox=\hbox{\insertsubsection}%
			\hbox{\ifdim\wd\beamer@tempbox>1pt{\hskip17.4pt\raise3pt\hbox{\vrule
						width0.4pt height7pt\vrule width 9pt
						height0.4pt}}\hskip1pt\hbox{\begin{minipage}[t]{7.5cm}\def\breakhere{}\insertsubsection\end{minipage}}\fi%
			}%
			\setbox\beamer@tempbox=\hbox{\insertshortframetitle}%
			\hbox{\ifdim\wd\beamer@tempbox>1pt{\hskip30.8pt\raise3pt\hbox{\vrule
						width0.4pt height7pt\vrule width 9pt
						height0.4pt}}\hskip1pt\hbox{\insertshortframetitle[width=7cm]}\fi%
			}%
			\vfil}%
	}%
	\ifbeamercolorempty[bg]{note page}{}{%
		\nointerlineskip%
		\insertvrule{.55\paperheight}{note page.bg}%
		\vskip-.55\paperheight%
	}%
	\vskip.25em
	\nointerlineskip
	\insertnote
}
\makeatother

\setbeamertemplate{note page}[mynotes]

\makeatletter
\setbeamertemplate{section page}{
	\centering
	\begin{minipage}{22em}
		\raggedright
		\usebeamercolor[fg]{section title}
		\usebeamerfont{section title}
		\thesection.~\insertsectionhead\\[-1ex]
		\usebeamertemplate*{progress bar in section page}
		\par
		\ifx\insertsubsectionhead\@empty\else%
		\usebeamercolor[fg]{subsection title}%
		\usebeamerfont{subsection title}%
		\thesection.\thesubsection~\insertsubsectionhead
		\fi
	\end{minipage}
	\par
	\vspace{\baselineskip}
}
\makeatother

\definecolor{myColor}{HTML}{4477AA}
\definecolor{pitchblack}{HTML}{000000}
\setbeamercolor{alerted text}{fg=myColor}
\setbeamercolor{frametitle}{bg=pitchblack, fg=white}

\makeatletter
\setlength{\metropolis@titleseparator@linewidth}{1pt}
\setlength{\metropolis@progressonsectionpage@linewidth}{2pt}
\setlength{\metropolis@progressinheadfoot@linewidth}{2pt}
\makeatother


% uncomment the following line if you want to hide the navigation symbols
%\beamertemplatenavigationsymbolsempty 

\title{Automaten und Formale Sprachen}
\subtitle{Theoretische Informatik 1}
\author{\texorpdfstring{Peter Kossek\newline\href{mailto:e0030414@ba-sachsen.de}{\scriptsize e0030414@ba-sachsen.de}}{Peter Kossek}}

\institute{Berufsakademie Sachsen, Staatliche Studienakademie Leipzig} % Deutsch
\date{}

\makeatletter
\@addtoreset{figure}{section}
\@addtoreset{table}{section}
\makeatother

\renewcommand{\thefigure}{\thesection.\arabic{figure}}

\renewcommand{\thetable}{\thesection.\arabic{table}}

\usepackage{nameref}
\makeatletter
\newcommand*{\currentname}{\@currentlabelname}
\makeatother

\usepackage{circuitikz}
\usepackage{tikz}
\usetikzlibrary{automata, positioning, arrows}

%strikeout text in math mode
\usepackage[normalem]{ulem}
\usepackage{amsmath}
\newcommand{\stkout}[1]{\ifmmode\text{\sout{\ensuremath{#1}}}\else\sout{#1}\fi}

\usepackage{tabularx} %same-width columns
\newcolumntype{Y}{>{\centering\arraybackslash}X} %columntype Y is centered column for tabularx

\setcounter{tocdepth}{1} % only sections in TOC

\usepackage{wasysym} %for lightning symbol

\usepackage{qrcode} %for QR codes

\begin{document}

\selectlanguage{ngerman}

\include{titel}

%\frame{
%	\frametitle{Sprachen im Kontext der Informatik}
%	
%	\begin{itemize}
%		\item \emph{Sprache:} Menge zulässiger Worte über einem Alphabet
%		\item \emph{Programm:} String über einem Alphabet von Eingabezeichen (Quellcode) bzw. binär über 0/1
%		\item \emph{Compiler:} 
%		\begin{itemize}
%			\item prüft die Zulässigkeit dieses Strings als Satz (Wort) einer Sprache
%			\item generiert Folge von Maschinenbefehlen
%		\end{itemize}
%		\item \emph{Lexikalische Analyse:} Übersetzer-Modul ("`Scanner"') fasst Symbolfolgen zu Schlüsselworten ("`Token"') zusammen
%		\item \emph{Syntaxanalyse:}
%		\begin{itemize}
%			\item "`Parser"' prüft, ob Folge von Token zulässig ist
%			\item Semantik-Prüfung erfolgt dabei nicht
%		\end{itemize}
%	\end{itemize}
%}

\include{logik}
\include{logischer_schluss}
\include{erfuellbarkeit}
\include{einfuehrung_sprachen_grammatiken}
\include{regulaere_sprachen_endliche_automaten}
\include{kellerautomaten}
\section{Kontextfreie Sprachen}

\begin{frame}{Typ-2 Grammatiken: Normalformen}
	\begin{itemize}
		\item Vereinfachungen für kontextfreie Grammatiken $G$
		\begin{itemize}
			\item $\varepsilon$-Regeln eliminieren: $S \rightarrow aAb; A \rightarrow aS|\varepsilon$\\
			wird zu	$S \rightarrow aAb|ab; A \rightarrow aS$
			\item $G$ reduzieren: nutzlose Variablen eliminieren
			\item Kettenregeln $A \rightarrow B$ entfernen
		\end{itemize}
		\item Chomsky-Normalform
		\begin{itemize}
			\item zu jeder $\varepsilon$-freien, reduzierten kontextfreien Grammatik $G$ ohne Kettenregeln lässt sich eine äquivalente Grammatik $G'$ angeben, die nur Regeln der Art $A \rightarrow a$ bzw. $A \rightarrow BC$ enthält.
			\item Konstuktion: Ausgangspunkt $A\rightarrow X_1X_2 \ldots X_m, m\geq 2, X_i \in V \cup \Sigma$
			\begin{itemize}
				\item für $X_i=a\in \Sigma$ neue Variable $X_a$ mit $X_a \rightarrow a$ einführen
				\item verbleibende Regeln $A \rightarrow B_1B_2\ldots B_m, B_i \in V, m\geq 3$:\\
				ersetze $A \rightarrow B_1D_1, D_1\rightarrow B_2D_2, \ldots, D_{m-2}\rightarrow B_{m-1}B_m$ ($D_i$ neue Variablen)
			\end{itemize}
		\end{itemize}
		\item Greibach-Normalform: Regeln vom Typ $A \rightarrow aB_1B_2 \ldots B_m|a$
	\end{itemize}
\end{frame}

\begin{frame}{Beispiel: Erzeugung Chomsky-Normalform}
	$G=(\{S, A, B, C, D, E\}, \{a, b\}, P, S)$\\
	$P=\{$\\
	\qquad\qquad $S \rightarrow bA \mid aB \mid C \mid Cb,$\\
	\qquad\qquad $A \rightarrow bAA \mid aS \mid a \mid bE \mid bDb,$\\
	\qquad\qquad $B \rightarrow aBB \mid bS \mid b,$\\
	\qquad\qquad $C \rightarrow Cc \mid c,$\\
	\qquad\qquad $D \rightarrow d \mid \varepsilon $\\
	\qquad$\}$
\end{frame}

\begin{frame}{Worterkennung: Cocke-Younger-Kasami-Algorithmus (CYK)}
	polynomiell auch für nichtdeterministische Sprachen
	\begin{itemize}
		\item Ziel: Worterkennung und Konstruktion Ableitungsbaum bottom-up
		\item Ausgangspunkt: $G$ in Chomsky-Normalform, $w=a_1a_2 \ldots a_n \in \Sigma^*$
		\item definiere $V[i,j]:=\{A \in V \mid A \underset{G}{\Rightarrow}^* w_{ij} \}$ mit\\
		$w_{ij}=a_ia_{i+1}\ldots a_{i+j-1}, |w_{ij}|=j, 1\leq i \leq n, 1\leq j \leq n+1-i$
		\item es gilt: $w\in L$ gdw. $S \in V[1,n]$
		\item schrittweise Berechnung (Rückführung auf kürzere Teilworte)\\
		$j=1: V[i,1]=\{X \in V \mid (X\rightarrow a_i) \in P\}$\\
		$j>1: V[i,j]=\bigcup_{l=1}^{j-1}\{A\in V\mid \exists (A\rightarrow BC): B \in V[i,l] \land C \in V[i+l, j-l] \}$
		\item Wortproblem mit polynomiellem Zeitaufwand $O(n^3)$ lösbar
		\item Grammatik ist mehrdeutig, falls ein $V[i,j]$ mehrere Einträge hat, die von $S$ aus erreicht werden können
	\end{itemize}
\end{frame}

\begin{frame}{Beispiel: CYK-Algorithmus}
	$$L=\{a^nb^n\mid n\geq 1\}$$
	Zugehörige Grammatik in Chomsky-Normalform:
	\begin{align*}
		G=(&\{A, B, C, S\}, \{a, b\}, P, S)\\
		P=\{&S \rightarrow AB,\\
		    &S \rightarrow AC,\\
		    &C\rightarrow SB,\\
		    &A\rightarrow a,\\
		    &B\rightarrow b\}
	\end{align*}
	
\end{frame}

\note{
	Betrachten des Wortes $w=aabb$\\
	\begin{columns}
		\column{.3\textwidth}
		\begin{itemize}
			\item[$(a)$] $V[1,1]=\{A\}$
			\item[$(a)$] $V[2,1]=\{A\}$
			\item[$(b)$] $V[3,1]=\{B\}$
			\item[$(b)$] $V[4,1]=\{B\}$
			\item[$(aa)$] $V[1,2]=\varnothing$
		\end{itemize}
		\column{.3\textwidth}
		\begin{itemize}
			\item[$(ab)$] $V[2,2]=\{S\}$
			\item[$(bb)$] $V[3,2]=\varnothing$
			\item[$(aab)$] $V[1,3]=\varnothing$
			\item[$(abb)$] $V[2,3]=\{C\} (l=2)$
			\item[$(aabb)$] $V[1,4]=\{S\} (l=1)$
		\end{itemize}
		\column{.4\textwidth}
		\centering
		\begin{tabular}{|c|c|c|c|}
			\hline
			$a$ & $a$ & $b$ & $b$ \\
			\hline
			\hline
			$\{A\}$ & $\{A\}$ & $\{B\}$ & $\{B\}$ \\
			\hline
			$\varnothing$ & $\{S\}$ & $\varnothing$ &  \\
			\hline
			$\varnothing$ & $\{C\}$ &  &  \\
			\hline
			$\{S\}$ &  &  &  \\
			\hline
		\end{tabular}
	\end{columns}
}

\begin{frame}{Pumping-Lemma für kontextfreie Sprachen}
	\begin{itemize}
		\item Pumping-Lemma: für jede kontextfreie Sprache $L$ gilt\\
		es existiert eine natürliche Zahl $n$ so, dass für alle $z \in L$ mit $|z|\geq n$ gilt:\\
		$z$ lässt sich zerlegen in $z=uvwxy$ mit
		\begin{itemize}
			\item $|vx|\geq 1$
			\item $|vwx| \leq n$
			\item $uv^iwx^iy \in L, i=0, 1, 2, \ldots$
		\end{itemize}
		\item Beweis:
		\begin{itemize}
			\item wähle $n=2^m$, $m$: Anzahl der Nichtterminalsymbole, wenn sich G in Chomsky-Normalform befindet
			\item mindestens eine Variable $A$ muss sich im Ableitungsbaum nach spätestens $m$ Schritten wiederholen;\\
			$A\underset{G}{\Rightarrow}^*\omega_1 A \omega_2$, wähle $v: \omega_1 \underset{G}{\Rightarrow}^* v, x: \omega_2 \underset{G}{\Rightarrow}^* x, w: A \underset{G}{\Rightarrow}^* w$;\\
			diese Ableitungen enthalten jeweils keine Wiederholungen von Variablen
		\end{itemize}
	\end{itemize}
\end{frame}

\begin{frame}{Beispiel: Pumping-Lemma für kontextfreie Sprachen}
	$L=\{a^ib^ic^i \mid i\geq 1 \}$
	\begin{itemize}
		\item $L$ ist kontext-sensitiv: Beweis Typ-1-Grammatik\\
		$G=(V, \Sigma, P, S), V=\{S, B, C\}$\\
		$P=\{S\rightarrow aSBC | aBC, CB \rightarrow BC,$\\
		\qquad $aB \rightarrow ab, bB \rightarrow bb, bC \rightarrow bc, cC \rightarrow cc \}$
		\item $L$ ist nicht kontextfrei: Beweis Pumping-Lemma\\
		wäre $L$ kontextfrei, sei $n$ das "`Pumping-$n$"'\\
		wähle $z=a^nb^nc^n \in L, |z|=3n \geq n$\\
		es ex. eine Zerlegung $z=uvwxy$ mit $|vx|\geq 1, |vwx| \leq n$\\
		$\rightarrow$ $vwx$ kann höchstens zwei verschiedene Buchstaben enthalten\\
		$\rightarrow$ in $uv^iwx^iy$ werden ein oder zwei Buchstaben "`aufgepumpt"', der dritte nicht\\
		$\rightarrow$ $uv^iwx^iy \notin L$ für $i \neq 1$ \lightning 
	\end{itemize}
\end{frame}

\begin{frame}{Abschlusseigenschaften kontextfreier Sprachen}
	\begin{itemize}
		\item Kontextfreie Sprachen sind abgeschlossen unter
		\begin{itemize}
			\item Vereinigung: $L_1\cup L_2$\qquad $S \rightarrow S_1|S_2$
			\item Konkatenation: $L_1L_2$ \qquad $S\rightarrow S_1S_2$
			\item Stern-Produkt: $L^*$ \qquad $S_0\rightarrow S|\varepsilon, S\rightarrow S_1|S_1S$
			\item Spiegelung: ersetze in Chomsky-Normalform $A\rightarrow BC$ durch $A \rightarrow CB$
		\end{itemize}
		\item Kontextfreie Sprachen sind \underline{nicht} abgeschlossen unter
		\begin{itemize}
			\item [a)] Durchschnitt\\
			Gegenbeispiel: $L_1=\{a^kb^kc^l\mid k,l\geq 0\}, L_2=\{a^kb^lc^l\mid k,l\geq 0\}$ sind kontextfrei\\
			$L=L_1\cap L_2$ ist nicht kontextfrei
			\item [b)] Komplement\\
			$L_1 \cap L_2 = \overline{\overline{L_1} \cup \overline{L_2}}$: Abgeschlossenheit unter Vereinigung und Komplement würde zu Abgeschlossenheit unter Durchschnitt führen\\
			$\rightarrow$ Widerspruch zu a) \lightning
			
		\end{itemize}
	\end{itemize}
\end{frame}

\begin{frame}{Abschlusseigenschaften deterministisch kontextfreier Sprachen}
	\begin{enumerate}[a)]
		\item Deterministisch kontextfreie Sprachen sind abgeschlossen unter Komplementbildung (beachte: Akzeptanz durch Endzustand!)
		\item Deterministisch kontextfreie Sprachen sind \underline{nicht} abgeschlossen unter
		\begin{enumerate}[i.]
			\item Durchschnitt $L_1 \cap L_2$
			\item Konkatenation $L_1L_2$
			\item Stern-Produkt $L^*$
			\item Vereinigung $L_1 \cup L_2$
		\end{enumerate}
		\item $L_1 \cap L_2$ ist allerdings det. kontextfrei (bzw. kontextfrei) falls $L_1$ regulär und $L_2$ det. kontextfrei (bzw. kontextfrei) ist
	\end{enumerate}
\end{frame}

\include{typ-0-sprachen_turing_maschinen}
% !TeX root = AuFS.tex
\section{Abschluss}

\begin{frame}{Empfehlung: JUG Saxony Day}
	\centering
	\includegraphics[width=0.7\linewidth]{images/JSD-Logo_gruen}
	\vspace{2em}
	\begin{itemize}
		\item September in Radebeul
		\item Freitickets für Studis
		\item Großartige Vorträge zu verschiedensten Themen
		\item Lecker Omnomnom
		\item \url{https://jugsaxony.org/day}
	\end{itemize}
\end{frame}

\begin{frame}{Kontakt}
	\begin{itemize}
		\item[{\includegraphics[width=3ex]{images/LI-In-Bug.png}}] \href{https://www.linkedin.com/in/peter-kossek-851b90195}{linkedin.com/in/peter-kossek-851b90195}
	
		\item[{\includegraphics[width=3ex]{images/XING_logo_neu.png}}] \href{https://www.xing.com/profile/Peter_Kossek2}{xing.com/profile/Peter\_Kossek2}
		
		\item[{\includegraphics[width=3ex]{images/telegram.png}}] \href{https://t.me/squinting_cyclops}{\texttt{@squinting\_cyclops}}
		
		\item[{\includegraphics[width=3ex]{images/instagram.png}}] \href{https://www.instagram.com/squinting.cyclops}{\texttt{@squinting.cyclops}}
	\end{itemize}

\end{frame}

\end{document}

